\chapter{1992年上海卷}
\section{填空题}

\begin{problem}
求值：\(\cos \frac{5\pi}{8}\cos \frac{\pi}{8}=\)\fillinblank{}.

\end{problem}

\begin{answer}
\(-\frac{\sqrt{2}}{4}\)
\end{answer}

\begin{problem}
\(\left(x + \frac{1}{x}\right)^8\) 的展开式中 \(\frac{1}{x^2}\) 的系数是\fillinblank{}（结果用数值表示）.

\end{problem}

\begin{answer}
56
\end{answer}

\begin{problem}
函数 \(y = \sin^2 x - \sin x \cos x + \cos^2 x\) 的最大值是\fillinblank{}.

\end{problem}

\begin{answer}
\(\frac{3}{2}\)
\end{answer}

\begin{problem}
方程 \(\log_5(x+1)-\log_{\frac{1}{5}}(x-3)=1\) 的解是\fillinblank{}.

\end{problem}

\begin{answer}
4
\end{answer}

\begin{problem}
计算 \((\sqrt{3} + \i)^6=\)\fillinblank{}（\(\i\) 是虚数单位）.

\end{problem}

\begin{answer}
\(-64\)
\end{answer}

\begin{problem}
如果直线 \(l\) 与直线 \(x + y - 1 = 0\) 关于 \(y\) 轴对称，那么直线 \(l\) 的方程是\fillinblank{}.

\end{problem}

\begin{answer}
\(x - y + 1 = 0\)
\end{answer}

\begin{problem}
已知圆台的下底面半径为 \(8\mathrm{cm}\)，高为 \(6\mathrm{cm}\)，母线与下底面成 \(45^\circ\) 的角，那么圆台的侧面积是\fillinblank{} \(\mathrm{cm}^2\)（结果中保留 \(\pi\)）.

\end{problem}

\begin{answer}
\(60\sqrt{2}\pi\)
\end{answer}

\begin{problem}
已知函数 \(y = f(x)\) 的反函数是 \(f^{-1}(x)=\sqrt{x}-1\)（\(x \geqslant 0\)），那么函数 \(y = f(x)\) 的定义域是\fillinblank{}.

\end{problem}

\begin{answer}
\([-1,+\infty)\)
\end{answer}

\begin{problem}
由 \(1\)，\(2\)，\(3\)，\(4\)，\(5\) 组成比 \(40000\) 小的没有重复数字的五位数的个数是\fillinblank{}（结果用数值表示）.

\end{problem}

\begin{answer}
72
\end{answer}

\begin{problem}
如图，直平行六面体 \(ABCD-A_1B_1C_1D_1\) 的上底面 \(ABCD\) 为菱形，\(\angle BAD = 60^\circ\)，侧面为正方形. \(E\)，\(F\) 分别是 \(A_1B_1\)，\(AA_1\) 的中点，\(M\) 是 \(AC\) 与 \(BD\) 的交点，则 \(EF\) 与 \(B_1M\) 所成角的大小为\fillinblank{}（用反三角函数表示）.

\bitmapfigure[max width=0.32\linewidth]{img/1992/shanghai/q10_fig1.png}

\end{problem}

\begin{answer}
\(\arcsin \frac{\sqrt{6}}{4}\)
\end{answer}

\section{单选题}

\begin{problem}
当 \(0 < a < 1\) 时，函数 \(y = a^x\) 和 \(y = (a - 1)x^2\) 的图象只可能是（\quad）

\choices
    {\choicebitmap[width=0.15\paperwidth]{img/1992/shanghai/q11_choiceA.png}}
    {\choicebitmap[width=0.15\paperwidth]{img/1992/shanghai/q11_choiceB.png}}
    {\choicebitmap[width=0.15\paperwidth]{img/1992/shanghai/q11_choiceC.png}}
    {\choicebitmap[width=0.15\paperwidth]{img/1992/shanghai/q11_choiceD.png}}

\end{problem}

\begin{answer}
D
\end{answer}

\begin{problem}
下列函数中，在区间 \((0,1]\) 上为减函数的是（\quad）

\choices
    {\(y = \log_a x\)}
    {\(y = x^a\)}
    {\(y = \cos x\)}
    {\(y = \arcsin x\)}

\end{problem}

\begin{answer}
C
\end{answer}

\begin{problem}
\(\alpha\)，\(\beta\) 是两个不重合的平面，在下列条件中，可判定平面 \(\alpha\) 与 \(\beta\) 平行的是（\quad）

\choices
    {\(\alpha\)，\(\beta\) 都垂直平面 \(\gamma\)}
    {\(\alpha\) 内不共线的三点到 \(\beta\) 的距离都相等}
    {\(l\)，\(m\) 是 \(\alpha\) 内两条直线，且 \(l \parallel \beta\)，\(m \parallel \beta\)}
    {\(l\)，\(m\) 是两条异面直线，且 \(l \parallel \alpha\)，\(m \parallel \alpha\)，\(l \parallel \beta\)，\(m \parallel \beta\)}

\end{problem}

\begin{answer}
D
\end{answer}

\begin{problem}
下列函数中，最小正周期为 \(\pi\) 的偶函数是（\quad）

\choices
    {\(y = \sin 2x\)}
    {\(y = \cos \frac{x}{2}\)}
    {\(y = \sin 2x + \cos 2x\)}
    {\(y = \frac{1 - \tan^2 x}{1 + \tan^2 x}\)}

\end{problem}

\begin{answer}
D
\end{answer}

\begin{problem}
在方程 \(\left\{\begin{array}{l} x = \sin \theta, \\ y = \cos 2\theta \end{array}\right.\)（\(\theta\) 为参数）所表示的曲线上的一个点的坐标是（\quad）

\choices
    {\((2,-7)\)}
    {\(\left(\frac{1}{3}, \frac{2}{3}\right)\)}
    {\(\left(\frac{1}{2}, \frac{1}{2}\right)\)}
    {\((1,0)\)}

\end{problem}

\begin{answer}
C
\end{answer}

\begin{problem}
抛物线 \((y + 2)^2 = 4(x + 1)\) 的焦点坐标是（\quad）

\choices
    {\((0,-2)\)}
    {\((2,2)\)}
    {\((1,-2)\)}
    {\((3,2)\)}

\end{problem}

\begin{answer}
A
\end{answer}

\begin{problem}
下列命题中的真命题是（\quad）

\choices
    {各侧面都是矩形的棱柱是长方体}
    {有两个相邻侧面是矩形的棱柱是直棱柱}
    {各侧面都是等腰三角形的四棱锥是正四棱锥}
    {有两个面互相平行，其余四个面都是等腰梯形的六面体是四棱台}

\end{problem}

\begin{answer}
B
\end{answer}

\begin{problem}
函数 \(y = \arccos \frac{1}{\sqrt{x}}\) 的值域是（\quad）

\choices
    {\(\left[0,\frac{\pi}{2}\right)\)}
    {\(\left(0,\frac{\pi}{2}\right]\)}
    {\(\left[0,\pi\right)\)}
    {\(\left(0,\pi\right]\)}

\end{problem}

\begin{answer}
A
\end{answer}

\begin{problem}
在极坐标系中，与圆 \(\rho = 4\sin \theta\) 相切的一条直线的方程是（\quad）

\choices
    {\(\rho \sin \theta = 2\)}
    {\(\rho \cos \theta = 2\)}
    {\(\rho \cos \theta = 4\)}
    {\(\rho \cos \theta = -4\)}

\end{problem}

\begin{answer}
B
\end{answer}

\begin{problem}
设集合 \(M = \{x \mid x > 2\}\)，\(P = \{x \mid x < 3\}\)，那么“\(x \in M\) 或 \(x \in P\)”是“\(x \in M \cap P\)”的（\quad）

\choices
    {充分条件但非必要条件}
    {必要条件但非充分条件}
    {充分必要条件}
    {非充分条件也非必要条件}

\end{problem}

\begin{answer}
B
\end{answer}

\section{解答题}

\begin{problem}
已知 \(\cos 2\alpha = \frac{7}{25}\)，\(\alpha \in \left(0, \frac{\pi}{2}\right)\)，\(\sin \beta = -\frac{5}{13}\)，\(\beta \in \left(\pi, \frac{3\pi}{2}\right)\)，求 \(\alpha + \beta\)（用反三角函数表示）.

\end{problem}

\begin{solution}
由 \(\cos 2\alpha = \frac{7}{25}\) 得 \(\cos \alpha = \sqrt{\frac{1 + \cos 2\alpha}{2}} = \frac{4}{5}\)，\(\sin \alpha = \frac{3}{5}\). 又由 \(\sin \beta = -\frac{5}{13}\) 且 \(\beta \in \left(\pi,\frac{3\pi}{2}\right)\)，得 \(\cos \beta = -\sqrt{1 - \sin^2 \beta} = -\frac{12}{13}\). 所以 \(\cos(\alpha + \beta) = \cos \alpha \cos \beta - \sin \alpha \sin \beta = -\frac{33}{65}\). 因为 \(\pi < \alpha + \beta < 2\pi\)，故 \(\alpha + \beta = \pi + \arccos \frac{33}{65}\).
\end{solution}

\begin{problem}
设 \(w = z + a\i\)，其中 \(a\) 是实数，\(z = \frac{(1 - 4\i)(1 + \i) + 2 + 4\i}{3 + 4\i}\)，且 \(|w| \leqslant \sqrt{2}\)，求 \(w\) 的辐角主值 \(\theta\) 的取值范围.

\end{problem}

\begin{solution}
先化简 \(z = \dfrac{(1 - 4\i)(1 + \i) + 2 + 4\i}{3 + 4\i}
= \dfrac{7 + \i}{3 + 4\i}
= 1 - \i\).
于是 \(w = z + a\i = 1 + (a - 1)\i\). 由 \(|w| \leqslant \sqrt{2}\) 得 \(\sqrt{1 + (a - 1)^2} \leqslant \sqrt{2}\)，所以 \(0 \leqslant a \leqslant 2\). 又 \(\tan\theta = a - 1\)，故 \(-1 \leqslant \tan\theta \leqslant 1\). 因此 \(\theta\) 的取值范围是 \(\left[0,\frac{\pi}{4}\right] \cup \left[\frac{7\pi}{4},2\pi\right)\).
\end{solution}

\begin{problem}
设动直线 \(l\) 垂直于 \(x\) 轴，且与椭圆 \(\frac{x^2}{4} + \frac{y^2}{2} = 1\) 交于 \(A\)，\(B\) 两点，\(P\) 是 \(l\) 上满足 \(|PA| \cdot |PB| = 1\) 的点，求点 \(P\) 的轨迹方程，并说明轨迹是什么图形.

\end{problem}

\begin{solution}
设点 \(P\) 的坐标为 \((x,y)\)，点 \(A\) 的坐标为 \((x,y_1)\)，则点 \(B\) 的坐标为 \((x,-y_1)\). 因为 \(A\)，\(B\) 两点在椭圆上，所以 \(\dfrac{x^2}{4} + \dfrac{y_1^2}{2} = 1\)（\circled{1}）.
又由 \(\circled{1}\) 得 \(-2 < x < 2\). 由题意，\(|PA| = |y_1 - y|\)，\(|PB| = |y_1 + y|\)，且 \(|PA|\cdot|PB| = 1\)，于是 \(|y_1 - y|\cdot|y_1 + y| = 1\)，即 \(y_1^2 - y^2 = \pm 1\). 若 \(y_1^2 = y^2 + 1\)，代入 \(\circled{1}\) 得 \(\frac{x^2}{2} + y^2 = 1\). 若 \(y_1^2 = y^2 - 1\)，代入 \(\circled{1}\) 得 \(\frac{x^2}{6} + \frac{y^2}{3} = 1\)（\(-2 < x < 2\)）. 因此，点 \(P\) 的轨迹是椭圆 \(\frac{x^2}{2} + y^2 = 1\)，以及椭圆 \(\frac{x^2}{6} + \frac{y^2}{3} = 1\) 夹在两直线 \(x = 2\) 与 \(x = -2\) 之间的部分.
\end{solution}

\begin{problem}
如图，圆锥的轴截面为等腰直角三角形 \(SAB\)，\(Q\) 为底面圆周上一点.

\begin{enumerate}
\item 如果 \(QB\) 的中点为 \(C\)，\(OH \perp SC\)，求证 \(OH \perp\) 平面 \(SBQ\)；
\item 如果 \(\angle AOQ = 60^{\circ}\)，\(QB = 2\sqrt{3}\)，求此圆锥的体积；
\item 如果二面角 \(A-SB-Q\) 的大小为 \(\arctan \frac{\sqrt{6}}{3}\)，求 \(\angle AOQ\) 的大小.
\end{enumerate}

\bitmapfigure[max width=0.58\linewidth]{img/1992/shanghai/q24_fig1.png}

\end{problem}

\begin{solution}
（1）连接 \(OC\). 因为 \(SQ = SB\)，\(OQ = OB\)，\(QC = CB\)，所以 \(QB \perp SC\)，且 \(QB \perp OC\). 故 \(QB \perp\) 平面 \(SOC\). 又因 \(OH \subset\) 平面 \(SOC\)，所以 \(QB \perp OH\). 再由 \(OH \perp SC\)，可得 \(OH \perp\) 平面 \(SBQ\).

（2）连接 \(AQ\). 因为 \(Q\) 为圆周上一点，\(AB\) 为直径，所以 \(AQ \perp QB\). 在 \(\mathrm{Rt}\triangle AQB\) 中，\(AB = \frac{2\sqrt{3}}{\cos 30^{\circ}} = 4\). 又因为 \(\triangle SAB\) 是等腰直角三角形，所以 \(SO = \frac{1}{2}AB = 2\). 故圆锥体积 \(V = \frac{1}{3}\pi \cdot OA^2 \cdot SO = \frac{8}{3}\pi\).

（3）过 \(Q\) 作 \(QM \perp AB\)，垂足为 \(M\). 因平面 \(SAB \perp\) 平面 \(ABQ\)，故 \(QM \perp\) 平面 \(SAB\). 再过 \(M\) 作 \(MP \perp SB\)，垂足为 \(P\)，连接 \(PQ\)，则由三垂线定理可得 \(QP \perp SB\). 因此 \(\angle MPQ\) 是二面角 \(A-SB-Q\) 的平面角，故 \(\angle MPQ = \arctan \frac{\sqrt{6}}{3}\). 设圆锥底面半径为 \(R\)，\(\angle AOQ = \alpha\). 在 \(\mathrm{Rt}\triangle OMQ\) 中，\(QM = R\sin \alpha\)，\(OM = R\cos \alpha\). 在 \(\mathrm{Rt}\triangle MPB\) 中，\(\angle PBM = 45^{\circ}\)，且 \(MB = R(1 + \cos \alpha)\)，故 \(MP = \frac{\sqrt{2}}{2}R(1 + \cos \alpha)\). 于是 \(\frac{R\sin \alpha}{\frac{\sqrt{2}}{2}R(1 + \cos \alpha)} = \frac{\sqrt{6}}{3}\)，即 \(\frac{\sin \alpha}{1 + \cos \alpha} = \frac{\sqrt{3}}{3}\). 因此 \(\tan\frac{\alpha}{2} = \frac{\sqrt{3}}{3}\)，所以 \(\angle AOQ = 60^{\circ}\).
\end{solution}

\begin{problem}
已知数列 \(\{a_n\}\)，\(a_n > 0\)（\(n \in \N\)），它的前 \(n\) 项的和记为 \(S_n\).

\begin{enumerate}
\item 如果 \(\{a_n\}\) 是一个首项为 \(a\)，公比为 \(q\)（\(0 < q \leqslant 1\)）的等比数列，且 \(G_n = a_1^2 + a_2^2 + \cdots + a_n^2\)（\(n \in \N\)），求 \(\lim\limits_{n \to \infty} \frac{S_n}{G_n}\)；
\item 如果 \(S_1^2\)，\(S_2^2\)，\(\cdots\)，\(S_n^2\)，\(\cdots\) 是一个首项为 3，公差为 1 的等差数列，试比较 \(S_n\) 与 \(3na_n\)（\(n \in \N\)）的大小.
\end{enumerate}

\end{problem}

\begin{solution}
（1）【法一】 当 \(q = 1\) 时，\(S_n = na\)，\(G_n = na^2\). 因为 \(a \neq 0\)，所以 \(\lim\limits_{n \to \infty} \frac{S_n}{G_n} = \frac{1}{a}\). 当 \(0 < q < 1\) 时，\(\lim\limits_{n \to \infty} S_n = \frac{a}{1 - q}\).
因为 \(\frac{a_n^2}{a_{n-1}^2} = q^2 < 1\)，所以 \(\{a_n^2\}\) 是首项为 \(a^2\)，公比为 \(q^2\) 的等比数列，\(\lim\limits_{n \to \infty} G_n = \frac{a^2}{1 - q^2}\).
所以 \(\lim\limits_{n \to \infty} \frac{S_n}{G_n} = \frac{\lim\limits_{n \to \infty} S_n}{\lim\limits_{n \to \infty} G_n} = \frac{1 + q}{a}\). 综上，当 \(q = 1\) 时，\(\lim\limits_{n \to \infty} \frac{S_n}{G_n} = \frac{1}{a}\)；当 \(0 < q < 1\) 时，\(\lim\limits_{n \to \infty} \frac{S_n}{G_n} = \frac{1 + q}{a}\).

【法二】 当 \(q = 1\) 时，同解一得 \(\lim\limits_{n \to \infty} \frac{S_n}{G_n} = \frac{1}{a}\).
当 \(0 < q < 1\) 时，\(S_n = \frac{a(1 - q^n)}{1 - q}\). 同解一得 \(\{a_n^2\}\) 是首项为 \(a^2\)，公比为 \(q^2\) 的等比数列，\(\lim\limits_{n \to \infty} G_n = \frac{a^2(1 - q^{2n})}{1 - q^2}\).
\(\lim\limits_{n \to \infty} \frac{S_n}{G_n} = \lim\limits_{n \to \infty} \frac{1 + q}{a(1 + q^n)} = \frac{1 + q}{a}\).
结论同解一.

（2）由题设知 \(S_1 = a_1 = \sqrt{3}\).
\(S_n^2 = 3 + (n - 1) = n + 2\)（\(n \in \N\)）.
因为 \(S_n > 0\)，所以 \(S_n = \sqrt{n + 2}\).
\(a_n = S_n - S_{n - 1} = \sqrt{n + 2} - \sqrt{n + 1}\)（\(n \geqslant 2\)）.
当 \(n = 1\) 时，\(S_1 < 3a_1\).
当 \(n \geqslant 2\) 时，
\(S_n - 3na_n = \sqrt{n + 2} - 3n\left(\sqrt{n + 2} - \sqrt{n + 1}\right) = 3n\sqrt{n + 1} - (3n - 1)\sqrt{n + 2}\).
由 \((3n\sqrt{n + 1})^2 - \left[(3n - 1)\sqrt{n + 2}\right]^2
= -3n^2 + 11n - 2
= -3\left(n - \frac{11 - \sqrt{97}}{6}\right)\left(n - \frac{11 + \sqrt{97}}{6}\right)\)，
可得当 \(n = 2\)，\(3\) 时，\(S_n > 3na_n\)；当 \(n \geqslant 4 (n \in \N)\) 时，\(S_n < 3na_n\).
综上所述，得
当 \(n = 1\) 及 \(n \geqslant 4 (n \in \N)\) 时，\(S_n < 3na_n\).
当 \(n = 2\)，\(3\) 时，\(S_n > 3na_n\).
\end{solution}

\begin{problem}
如图，已知双曲线 \(C: (1 - a^2)x^2 + a^2y^2 - a^2 = 0\)（参数 \(a > 0\)）. 若 \(C\) 的上半支的顶点为 \(A\)，且与直线 \(y = -x\) 交于点 \(P\). 以 \(A\) 为焦点，\(M(0,m)\) 为顶点的开口向下的抛物线通过点 \(P\). 当 \(C\) 的一条渐近线的斜率在区间 \(\left[\frac{\sqrt{3}}{2}, \frac{2\sqrt{2}}{3}\right]\) 上变化时，求直线 \(PM\) 斜率的最大值.

\bitmapfigure[max width=0.38\linewidth]{img/1992/shanghai/q26_fig1.png}

\end{problem}

\begin{solution}
【法一】 因为 \((1 - a^2)x^2 + a^2y^2 - a^2 = 0\) 是双曲线，所以 \(1 - a^2 < 0\)，从而它可化为 \(y^2 - \dfrac{x^2}{a^2 - 1} = 1\).
双曲线的两条渐近线的斜率为 \(\pm \frac{\sqrt{a^2 - 1}}{a}\). 于是 \(\frac{\sqrt{3}}{2} \leqslant \frac{\sqrt{a^2 - 1}}{a} \leqslant \frac{2\sqrt{2}}{3}\)，即 \(4 \leqslant a^2 \leqslant 9\)，又由 \(a > 0\) 得 \(a\) 的取值范围是 \([2,3]\).
把 \(y = -x\) 代入双曲线方程，得点 \(P\) 的坐标为 \((-a,a)\).
双曲线上半支的顶点为 \(A(0,1)\).
以 \(A(0,1)\) 为焦点，\(M(0,m)\) 为顶点的开口向下的抛物线方程为 \(x^2 = -4(m - 1)(y - m)\).
因为 \(P(-a,a)\) 在抛物线上，所以 \(a^2 = -4(m - 1)(a - m)\)（ \circled{1} ）.
即 \(m^2 - (1 + a)m + \left(a - \frac{a^2}{4}\right) = 0\). 解得 \(m = \dfrac{1}{2}\left(1 + a \pm \sqrt{1 - 2a + 2a^2}\right)\)，
根据题设抛物线的顶点位置，由 \circled{1} 可知 \(m > 1\)，\(m > a\)，所以 \(m > \dfrac{1 + a}{2}\).
因此只取 \(m = \dfrac{1}{2}\left(1 + a + \sqrt{1 - 2a + 2a^2}\right)\).
直线 \(PM\) 的斜率为 \(k = \dfrac{m - a}{0 - (-a)} = \dfrac{1}{2a}\left(1 - a + \sqrt{1 - 2a + 2a^2}\right)\)（ \circled{2} ）.
在 \circled{2} 中，令 \(u = \dfrac{1}{a}\)，有 \(k = \dfrac{1}{2}\left(u - 1 + \sqrt{u^2 - 2u + 2}\right) = \dfrac{1}{2}\left[\sqrt{(1 - u)^2 + 1} - (1 - u)\right]\).
在 \([2,3]\) 上，\(u\) 是减函数，\(1 - u\) 是增函数.由 \(1 - u > 0\)，可知 \((1 - u)^3\) 也是增函数，从而 \(k\) 在 \([2,3]\) 上是减函数. 故当 \(a = 2\) 时，\(k\) 取得最大值 \(k_{\max} = \dfrac{\sqrt{5} - 1}{4}\).

【法二】 同解一得双曲线方程为 \(y^2 - \dfrac{x^2}{a^2 - 1} = 1\).
\(a\) 的取值范围为 \([2,3]\)，
点 \(P\) 的坐标为 \((-a,a)\).
以 \(A\) 为焦点，\(M\) 为顶点的开口向下的抛物线方程为 \(x^2 = -4(m - 1)(y - m)\).
\(a^2 = -4(m - 1)(a - m)\)（ \circled{1} ）.
\(k = \dfrac{m - a}{a}\)（ \circled{2} ）.
由 \circled{2} 得 \(m = ka + a\)，代入 \circled{1} 得 \(a^2 = 4(ka + a - 1)ka\).
因为 \(a \neq 0\)，所以 \((4k^2 + 4k - 1)a = 4k\).
根据题设抛物线顶点的位置，由 \circled{1} 可知 \(m > a\)，所以 \(k > 0\). 因此 \(4k^2 + 4k - 1 > 0\).
这样 \(a = 4k/(4k^2 + 4k - 1)\).
由 \(a\) 的取值范围，得 \(2 \leqslant \dfrac{4k}{4k^2 + 4k - 1} \leqslant 3\)（ \circled{3} ）.
由 \(4k^2 + 4k - 1 > 0\)，得 \(k > \dfrac{\sqrt{2} - 1}{2}\). 于是，可由 \circled{3} 解得 \(\dfrac{\sqrt{13} - 2}{6} \leqslant k \leqslant \dfrac{\sqrt{5} - 1}{4}\).
因此 \(k_{\max} = \dfrac{\sqrt{5} - 1}{4}\).
\end{solution}
